\scanpage{17}
\mspage{16}

définit sur $X$ une structure de \deleted{\uline{boucle topologique}} \uline{maille}, \deleted{dont $A$ est le} \uline{noeud} $A$, si \deleted{$(\widetilde X,\{\widetilde A_1,\widetilde A_2\})$} \deleted{est} $A$ est régulièrement immergé dans $X$ (alors $\widetilde A_{\deleted{\unsure{2}}}$ est régulièrement immergé dans $\widetilde X$, donc $\widetilde A_1$, $\widetilde A_2$ aussi) et si $(\widetilde X,\{\widetilde A_1,\widetilde A_2\})$ satisfait les conditions a) b) c) des ficelles top. (donc est une ficelle top. si $\widetilde A$ est rég. immergé dans $\widetilde X$).
\margin{\unsure{réciproq.}}

Considérons
\[
\widetilde X\amalg_{\widetilde A_1}(\widetilde A_1\longrightarrow A)
\amalg_{\widetilde A_2}(\widetilde A_2\longrightarrow A),
\]
on trouve un espace $\overline X=\operatorname{Dec}_{\beta}(A,X)$, avec deux parties (indexées par $\beta$). can. isom. à $A$, $\{A_\rho,A_{\rho'}\}$, \deleted{\unsure{A}} i.e. une partie isom. à $A\times\beta$
\margin{$\beta=\{\rho,\rho'\}$ la birive}
\diagram{\begin{tikzcd}[ampersand replacement=\&]
\overline X\arrow[r]\&X
\end{tikzcd}}
induisant un isom
\diagram{\begin{tikzcd}[ampersand replacement=\&]
\overline X\setminus\{\underbrace{A_\rho\cup A_{\rho'}}_{A\times\unsure{\rho|\beta}}\}\arrow[r]\&X\setminus A
\end{tikzcd}}
et des homéom
\[
A_\rho\overset{\sim}{\longrightarrow} A,\qquad A_{\rho'}\simeq A
\]
i.e. la projection can.
\diagram{\begin{tikzcd}[ampersand replacement=\&]
A\times\beta\arrow[r,dash]\&A
\end{tikzcd}}

\rotatebox{90}{\reflectbox{\unsure{EXTRA STRONG}}}
\rotatebox{90}{\reflectbox{\unsure{EXTRA STRONG}}}
\mspage{16}
